
%---------------------------------------------------------------
% The "*" following chapter or section commands omits chapter/
% section numbers.  It also does not include the chapter/section
% in the table of contents -- the \addcontentsline can be used
% to manually force its entry.
%---------------------------------------------------------------

\chapter*{Introduction}
\addcontentsline{toc}{chapter}{Introduction}

The continued reduction of the transistor feature size has led to the increase in the vulnerability of modern circuits to error. This effect, in turn, has made modern circuits more susceptible to radiation induced errors in space and terrestrial environments. A radiation induced error, commonly referred to as a soft error, occurs when a high energy particle from space or packaging strikes a transistor. Shown in Fig. \ref{strike}, the particle deposits energy in the active volume generating electron-hole pairs. This mechanism creates a new diffusion region that could allow a non-conducting device to temporarily conduct current. The mechanism causes a temporary voltage pulse, referred to as a single event transient (SET) to occur at the device.

\begin{figure}[!htbp]
	\centering
	\includegraphics[width=0.45\linewidth]{Figures/StrikeFig}
	%where an .eps filename suffix will be assumed under latex, 
	%and a .pdf suffix will be assumed for pdflatex; or what has been declared
	%via \DeclareGraphicsExtensions.
	\caption{A energetic particle creating electron-hole pairs in a transistor.}
	\label{strike}
\end{figure} 

In the case of a radiation particle striking a memory storage element, the data stored in the device can be changed. If the device is loses its value by a single particle, this is referred to as a single event upset (SEU). Additionally, due to the constant scaling down of the feature size, a single particle can also strike two transistors simultaneously which is referred to as a double node upset (DNU).